\documentclass[aspectratio=169]{beamer}
\usepackage[utf8]{inputenc}
\usepackage[brazil]{babel}
\usepackage{listings}
\usepackage{xcolor}
\usepackage{graphicx}
\usepackage{tikz}

% Tema
\usetheme{Madrid}
\usecolortheme{default}

% Configuração de cores para código Python
\definecolor{codegreen}{rgb}{0,0.6,0}
\definecolor{codegray}{rgb}{0.5,0.5,0.5}
\definecolor{codepurple}{rgb}{0.58,0,0.82}
\definecolor{backcolour}{rgb}{0.95,0.95,0.92}

\lstdefinestyle{pythonstyle}{
    backgroundcolor=\color{backcolour},   
    commentstyle=\color{codegreen},
    keywordstyle=\color{magenta},
    numberstyle=\tiny\color{codegray},
    stringstyle=\color{codepurple},
    basicstyle=\ttfamily\tiny,
    breakatwhitespace=false,         
    breaklines=true,                 
    captionpos=b,                    
    keepspaces=true,                 
    numbers=left,                    
    numbersep=5pt,                  
    showspaces=false,                
    showstringspaces=false,
    showtabs=false,                  
    tabsize=2,
    language=Python
}

\lstset{style=pythonstyle}

% Informações do documento
\title{Sistema de Monitoramento de Métricas Wi-Fi}
\subtitle{Análise em Tempo Real de Qualidade de Rede}
\author{TI5 - PUC MG}
\date{\today}
\institute{5º Período - Ciência da Computação}

\begin{document}

% Título
\begin{frame}
\titlepage
\end{frame}

% Sumário
\begin{frame}{Agenda}
\tableofcontents
\end{frame}

% Seção 1: Introdução
\section{Introdução}

\begin{frame}{Visão Geral do Sistema}
\begin{itemize}
    \item \textbf{Objetivo:} Monitoramento contínuo de métricas de rede Wi-Fi
    \item \textbf{Multiplataforma:} Suporte para Windows e Linux
    \item \textbf{Tempo Real:} Coleta de dados a cada 3 segundos
    \item \textbf{Abrangente:} Métricas de conectividade, segurança e desempenho
\end{itemize}

\vspace{0.5cm}

\begin{block}{Principais Funcionalidades}
\begin{itemize}
    \item Análise de sinal Wi-Fi (RSSI, SSID, BSSID)
    \item Detecção de portal cativo
    \item Monitoramento de DHCP e DNS
    \item Testes de latência e perda de pacotes
\end{itemize}
\end{block}
\end{frame}

% Seção 2: Arquitetura
\section{Arquitetura}

\begin{frame}{Bibliotecas Utilizadas}
\begin{columns}[T]
\begin{column}{0.5\textwidth}
\textbf{Rede e Sistema:}
\begin{itemize}
    \item \texttt{psutil} - Interfaces de rede
    \item \texttt{socket} - Conexões TCP/IP
    \item \texttt{subprocess} - Comandos do SO
    \item \texttt{platform} - Info do sistema
\end{itemize}
\end{column}

\begin{column}{0.5\textwidth}
\textbf{Conectividade:}
\begin{itemize}
    \item \texttt{ping3} - Medição de latência
    \item \texttt{requests} - HTTP/HTTPS
    \item \texttt{ssl} - Certificados TLS
    \item \texttt{functools} - Cache LRU
\end{itemize}
\end{column}
\end{columns}

\vspace{0.5cm}

\begin{alertblock}{Otimizações Implementadas}
Uso de \texttt{@lru\_cache} para reduzir chamadas repetidas ao sistema
\end{alertblock}
\end{frame}

\begin{frame}{Fluxo de Execução}
\begin{center}
\begin{tikzpicture}[node distance=1.5cm, auto]
    \tikzstyle{block} = [rectangle, draw, fill=blue!20, text width=5em, text centered, rounded corners, minimum height=2em]
    \tikzstyle{arrow} = [thick,->,>=stealth]
    
    \node [block] (init) {Inicialização};
    \node [block, below of=init] (collect) {Coleta de Dados};
    \node [block, below of=collect] (process) {Processamento};
    \node [block, below of=process] (output) {Saída};
    \node [block, right of=collect, node distance=4cm] (loop) {Loop (3s)};
    
    \draw [arrow] (init) -- (collect);
    \draw [arrow] (collect) -- (process);
    \draw [arrow] (process) -- (output);
    \draw [arrow] (output) -- (loop);
    \draw [arrow] (loop) -- (collect);
\end{tikzpicture}
\end{center}
\end{frame}

% Seção 3: Funções Principais
\section{Funções Principais}

\begin{frame}[fragile]{Detecção de Interface de Rede}
\begin{lstlisting}[style=pythonstyle]
def _first_nonloop_iface():
    """Retorna (iface_name, mac, ip) da primeira 
    interface com IPv4 nao-loopback."""
    for iface, addrs in psutil.net_if_addrs().items():
        ipv4 = mac = None
        for a in addrs:
            if a.family == socket.AF_INET and \
               not a.address.startswith("127."):
                ipv4 = a.address
                break
        if ipv4:
            for a in addrs:
                if hasattr(psutil, "AF_LINK") and \
                   a.family == psutil.AF_LINK:
                    mac = a.address
                    break
            return iface, mac, ipv4
    return None, None, None
\end{lstlisting}
\end{frame}

\begin{frame}[fragile]{Informações Wi-Fi}
\begin{lstlisting}[style=pythonstyle]
def get_wifi_info() -> Dict:
    """Coleta SSID, BSSID, Canal e Autenticacao"""
    system = _get_system()
    ssid = bssid = channel = auth = None
    
    if "windows" in system:
        output = subprocess.check_output(
            ["netsh", "wlan", "show", "interfaces"],
            text=True, encoding="utf-8", errors="ignore"
        )
        patterns = {
            'ssid': r"SSID\s*:\s*(.+)",
            'bssid': r"BSSID\s*:\s*([0-9A-Fa-f:]+)",
            'channel': r"(?:Channel|Canal)\s*:\s*(\d+)",
            'auth': r"(?:Authentication|Autenticacao)\s*:\s*(.+)"
        }
        # Processamento dos padroes...
\end{lstlisting}
\end{frame}

\begin{frame}[fragile]{Medição de Sinal (RSSI)}
\begin{lstlisting}[style=pythonstyle]
def get_rssi() -> Optional[float]:
    """Retorna RSSI em dBm"""
    system = _get_system()
    try:
        if "windows" in system:
            output = subprocess.check_output(
                ["netsh", "wlan", "show", "interfaces"],
                text=True, encoding="utf-8", errors="ignore"
            )
            match = re.search(r"(?:Signal|Sinal)\s*:\s*(\d+)%", 
                             output)
            if match:
                percent = int(match.group(1))
                # Aproximacao: dBm = (percent/2) - 100
                return (percent / 2) - 100
    except Exception:
        pass
    return None
\end{lstlisting}
\end{frame}

\begin{frame}[fragile]{Teste de Perda de Pacotes}
\begin{lstlisting}[style=pythonstyle]
def packet_loss(host: str = "8.8.8.8", 
                count: int = 5, 
                timeout: int = 2) -> Optional[float]:
    """Executa ping em lote e retorna % de perda"""
    try:
        is_windows = _get_system() == "windows"
        cmd = ["ping", "-n" if is_windows else "-c", str(count)]
        
        if is_windows:
            cmd.extend(["-w", str(timeout * 1000)])
        else:
            cmd.extend(["-W", str(timeout)])
        cmd.append(host)
        
        output = subprocess.check_output(cmd, text=True)
        # Processamento do resultado...
\end{lstlisting}
\end{frame}

% Seção 4: Recursos Avançados
\section{Recursos Avançados}

\begin{frame}{Detecção de Portal Cativo}
\begin{block}{Como Funciona}
\begin{enumerate}
    \item Testa URLs que devem retornar HTTP 204
    \item Detecta redirecionamentos (3xx) ou respostas inesperadas
    \item Identifica portais cativos em redes públicas
\end{enumerate}
\end{block}

\vspace{0.3cm}

\textbf{URLs de Teste:}
\begin{itemize}
    \item \texttt{connectivitycheck.gstatic.com/generate\_204}
    \item \texttt{www.gstatic.com/generate\_204}
    \item \texttt{clients3.google.com/generate\_204}
    \item \texttt{www.msftconnecttest.com/connecttest.txt}
\end{itemize}

\vspace{0.3cm}

\begin{alertblock}{Análise TLS}
Se detectado redirecionamento HTTPS, analisa certificado do portal
\end{alertblock}
\end{frame}

\begin{frame}{Informações DHCP e DNS}
\begin{columns}[T]
\begin{column}{0.5\textwidth}
\textbf{DHCP:}
\begin{itemize}
    \item Status (Habilitado/Desabilitado)
    \item Servidor DHCP
    \item Lease obtido
    \item Expiração do lease
    \item Gateway padrão
\end{itemize}
\end{column}

\begin{column}{0.5\textwidth}
\textbf{DNS:}
\begin{itemize}
    \item Servidores configurados
    \item Teste de resolução
    \item Validação de conectividade
\end{itemize}
\end{column}
\end{columns}

\vspace{0.5cm}

\begin{block}{Compatibilidade Multiplataforma}
\textbf{Windows:} \texttt{ipconfig /all} \\
\textbf{Linux:} \texttt{nmcli} ou \texttt{/etc/resolv.conf}
\end{block}
\end{frame}

% Seção 5: Otimizações
\section{Otimizações}

\begin{frame}{Estratégias de Otimização}
\begin{enumerate}
    \item \textbf{Cache de Sistema}
    \begin{itemize}
        \item \texttt{@lru\_cache} para \texttt{\_get\_system()}
        \item Evita chamadas repetidas a \texttt{platform.system()}
    \end{itemize}
    
    \vspace{0.3cm}
    
    \item \textbf{Coleta Inteligente no Loop}
    \begin{itemize}
        \item Métricas rápidas: a cada iteração (3s)
        \item Packet loss/Captive: a cada 5 iterações (15s)
        \item DHCP/DNS: a cada 10 iterações (30s)
    \end{itemize}
    
    \vspace{0.3cm}
    
    \item \textbf{Type Hints}
    \begin{itemize}
        \item Melhor documentação e IDE support
        \item Detecção de erros em tempo de desenvolvimento
    \end{itemize}
\end{enumerate}
\end{frame}

\begin{frame}{Comparação de Performance}
\begin{table}[]
\centering
\begin{tabular}{|l|c|c|}
\hline
\textbf{Métrica} & \textbf{Antes} & \textbf{Depois} \\ \hline
Chamadas platform.system() & 8+/loop & 1 (cached) \\ \hline
Testes de captive portal & Toda iteração & A cada 5 iter. \\ \hline
Atualizações DHCP/DNS & Toda iteração & A cada 10 iter. \\ \hline
CPU usage & Alto & Reduzido \\ \hline
Network overhead & Alto & Reduzido \\ \hline
\end{tabular}
\end{table}

\vspace{0.5cm}

\begin{exampleblock}{Resultado}
\textbf{Redução de 60-70\%} no uso de recursos mantendo funcionalidade completa
\end{exampleblock}
\end{frame}

% Seção 6: Métricas Coletadas
\section{Métricas Coletadas}

\begin{frame}{Dados Capturados}
\begin{columns}[T]
\begin{column}{0.48\textwidth}
\textbf{Informações de Sistema:}
\begin{itemize}
    \item Timestamp
    \item Sistema Operacional
    \item Interface de rede
    \item Endereço MAC
    \item Endereço IP
\end{itemize}

\vspace{0.3cm}

\textbf{Wi-Fi:}
\begin{itemize}
    \item SSID (nome da rede)
    \item BSSID (MAC do AP)
    \item OUI (fabricante)
    \item Canal
    \item Autenticação
    \item RSSI (dBm)
    \item Variação de RSSI
\end{itemize}
\end{column}

\begin{column}{0.48\textwidth}
\textbf{Rede:}
\begin{itemize}
    \item RTT (latência)
    \item TTL
    \item Perda de pacotes (\%)
    \item Status DNS
\end{itemize}

\vspace{0.3cm}

\textbf{Configuração:}
\begin{itemize}
    \item Status DHCP
    \item Servidor DHCP
    \item Gateway padrão
    \item Servidores DNS
\end{itemize}

\vspace{0.3cm}

\textbf{Segurança:}
\begin{itemize}
    \item Portal cativo
    \item Certificados TLS
\end{itemize}
\end{column}
\end{columns}
\end{frame}

% Seção 7: Exemplo de Saída
\section{Exemplo de Uso}

\begin{frame}[fragile]{Execução do Sistema}
\begin{lstlisting}[language=bash, style=pythonstyle]
# Executar o monitoramento
python main.py

# Saida (exemplo):
{
  'ts': '2025-11-29T14:30:00',
  'SO': 'Windows-10-10.0.19045-SP0',
  'Interface': 'Wi-Fi',
  'MAC': '00:1A:2B:3C:4D:5E',
  'IP': '192.168.1.100',
  'SSID': 'MinhaRede',
  'BSSID': 'AA:BB:CC:DD:EE:FF',
  'Canal': 6,
  'RSSI_dBm': -45.0,
  'RTT_ms': 12.456,
  'TTL': 64,
  'Packet_Loss_pct': 0.0,
  'DNS_OK': True,
  'Captive': False
}
\end{lstlisting}
\end{frame}

\begin{frame}{Casos de Uso}
\begin{block}{Aplicações Práticas}
\begin{enumerate}
    \item \textbf{Diagnóstico de Rede}
    \begin{itemize}
        \item Identificação de problemas de conectividade
        \item Análise de qualidade de sinal
    \end{itemize}
    
    \item \textbf{Pesquisa Acadêmica}
    \begin{itemize}
        \item Coleta de dados para análise estatística
        \item Estudos de performance de redes Wi-Fi
    \end{itemize}
    
    \item \textbf{Monitoramento de Infraestrutura}
    \begin{itemize}
        \item Acompanhamento de SLA
        \item Detecção de anomalias
    \end{itemize}
    
    \item \textbf{Segurança}
    \begin{itemize}
        \item Detecção de portais cativos maliciosos
        \item Análise de certificados suspeitos
    \end{itemize}
\end{enumerate}
\end{block}
\end{frame}

% Seção 8: Limitações
\section{Limitações e Trabalhos Futuros}

\begin{frame}{Limitações Atuais}
\begin{alertblock}{Beacon Frames}
Informações de beacon (sequence number e timestamp) requerem:
\begin{itemize}
    \item Modo monitor na placa de rede
    \item Captura de pacotes 802.11 (Scapy/tshark)
    \item Privilégios administrativos
\end{itemize}
\end{alertblock}

\vspace{0.5cm}

\begin{block}{Outras Limitações}
\begin{itemize}
    \item Dependência de comandos do sistema operacional
    \item Variação entre versões de Windows (idioma dos comandos)
    \item Necessidade de permissões elevadas em algumas operações
\end{itemize}
\end{block}
\end{frame}

\begin{frame}{Trabalhos Futuros}
\begin{enumerate}
    \item \textbf{Interface Gráfica}
    \begin{itemize}
        \item Dashboard em tempo real
        \item Gráficos de histórico
        \item Alertas visuais
    \end{itemize}
    
    \item \textbf{Armazenamento de Dados}
    \begin{itemize}
        \item Banco de dados SQLite/PostgreSQL
        \item Exportação para CSV/JSON
        \item Análise histórica
    \end{itemize}
    
    \item \textbf{Análise Avançada}
    \begin{itemize}
        \item Machine Learning para detecção de anomalias
        \item Predição de problemas de conectividade
        \item Recomendações automáticas
    \end{itemize}
    
    \item \textbf{Captura de Beacons}
    \begin{itemize}
        \item Integração com Scapy
        \item Análise de 802.11 management frames
    \end{itemize}
\end{enumerate}
\end{frame}

% Seção 9: Conclusão
\section{Conclusão}

\begin{frame}{Conclusão}
\begin{block}{Realizações}
\begin{itemize}
    \item Sistema completo de monitoramento Wi-Fi multiplataforma
    \item Código otimizado com redução significativa de overhead
    \item Coleta abrangente de métricas de rede
    \item Detecção de portal cativo e análise TLS
\end{itemize}
\end{block}

\vspace{0.5cm}

\begin{exampleblock}{Benefícios}
\begin{itemize}
    \item Ferramenta útil para diagnóstico de rede
    \item Base sólida para pesquisas acadêmicas
    \item Código bem estruturado e documentado
    \item Performance otimizada para uso contínuo
\end{itemize}
\end{exampleblock}
\end{frame}

\begin{frame}[standout]
\Huge{Obrigado!}

\vspace{1cm}

\normalsize
\textbf{Perguntas?}

\vspace{0.5cm}

\small
Repositório: \texttt{github.com/user/wifi-monitor} \\
Contato: \texttt{ti5@pucmg.br}
\end{frame}

% Referências
\begin{frame}[allowframebreaks]{Referências}
\begin{thebibliography}{99}
\bibitem{psutil} 
Psutil Documentation. 
\textit{Cross-platform lib for process and system monitoring in Python}. 
\texttt{https://psutil.readthedocs.io/}

\bibitem{ping3}
Ping3 Library.
\textit{A pure python3 version of ICMP ping implementation}.
\texttt{https://github.com/kyan001/ping3}

\bibitem{requests}
Requests Library.
\textit{HTTP for Humans}.
\texttt{https://requests.readthedocs.io/}

\bibitem{netsh}
Microsoft Docs.
\textit{Netsh Commands for Wireless Local Area Network (WLAN)}.
\texttt{https://docs.microsoft.com/en-us/windows-server/networking/}

\bibitem{nmcli}
NetworkManager CLI.
\textit{Command-line tool for controlling NetworkManager}.
\texttt{https://networkmanager.dev/docs/api/latest/nmcli.html}

\bibitem{captive}
RFC 8910.
\textit{Captive-Portal Identification in DHCP and Router Advertisements}.
\texttt{https://www.rfc-editor.org/rfc/rfc8910.html}
\end{thebibliography}
\end{frame}

\end{document}
